% GENERATED FILE. Edit canonical lessons, then run npm run generate:books.
% canonical-source-hash: fnv1a64:7f90e3dd045cb2f5
% canonical-lessons: KA-C535-vadya, KA-C535-vine, KA-C535-kolalu, KA-C535-pitilu, KA-C535-barisu

\chapter{Instruments}
\label{ch:ka-instruments}
\hlchaptermodality{535}

\begin{chapteropening}
\textbf{By the end of this chapter:} \emph{I can say a musical instrument, a veena, a flute, a violin and to beat a drum.}

Complete the last lesson of chapter 535: I can say a musical instrument, a veena, a flute, a violin and to beat a drum.
\end{chapteropening}

\section[\texorpdfstring{\kn{ವಾದ್ಯ} (vādya)}{ವಾದ್ಯ (vādya)}]{\kn{ವಾದ್ಯ} (vādya) — a musical instrument}

\label{lesson:KA-C535-vadya}



\begin{quote}
Before the new one: say the Kannada for to hum a tune, then the Kannada for to play an instrument.
\end{quote}

\subsection*{You'll want to know: \kn{ವಾದ್ಯ}}
\textbf{\kn{ವಾದ್ಯ}} — \emph{vādya} — \textquotedblleft{}a musical instrument\textquotedblright{}.

\begin{quote}
\kn{ಮದುವೆಯಲ್ಲಿ} \kn{ವಾದ್ಯಗಳ} \kn{ಸದ್ದು} \kn{ಜೋರಾಗಿತ್ತು}.
\end{quote}

\begin{quote}
\emph{maduveyalli vādyagaḷa saddu jōrāgittu.}
\end{quote}

\begin{quote}
At the wedding the instruments were loud.
\end{quote}

Plural: \kn{ವಾದ್ಯಗಳು}, \emph{vādyagaḷu}.

\kn{ವಾದ್ಯಗಾರ}, a player.

\begin{cousinweb}
\textbf{\kn{ವಾದ್ಯ}} (\emph{vādya}) is borrowed from Sanskrit \emph{vādya}, \textquotedblleft{}to be played; music, a musical instrument\textquotedblright{} (CDIAL 11511).
\end{cousinweb}

\begin{grammarlens}[title={What you've built}]
One more word for describing what you lived and felt.
\end{grammarlens}

\subsection*{Your turn}
\begin{itemize}
\raggedright
  \item \emph{Say it:} \emph{vādya}
  \item \emph{Say it:} \emph{vādya}, once more
  \item \emph{Say it:} \emph{maduveyalli vādyagaḷa saddu jōrāgittu.} — At the wedding the instruments were loud.
  \item \emph{Recall:} say the Kannada for to hum a tune, then the Kannada for to play an instrument, then say \emph{vādya} again
\end{itemize}

\begin{tcolorbox}[breakable,colback=teal!4,colframe=teal!35!black,title={Before you move on}]
What does \kn{ವಾದ್ಯ} mean? (\textquotedblleft{}A musical instrument\textquotedblright{}.) Say it once more.
\end{tcolorbox}
\section[\texorpdfstring{\kn{ವೀಣೆ} (vīṇe)}{ವೀಣೆ (vīṇe)}]{\kn{ವೀಣೆ} (vīṇe) — a veena}

\label{lesson:KA-C535-vine}



\begin{quote}
Before the new one: say the Kannada for to play an instrument, then the Kannada for a musical instrument.
\end{quote}

\subsection*{You'll want to know: \kn{ವೀಣೆ}}
\textbf{\kn{ವೀಣೆ}} — \emph{vīṇe} — \textquotedblleft{}a veena\textquotedblright{}.

\begin{quote}
\kn{ವೀಣೆಯ} \kn{ನಾದ} \kn{ತುಂಬಾ} \kn{ಮಧುರ}.
\end{quote}

\begin{quote}
\emph{vīṇeya nāda tumbā madhura.}
\end{quote}

\begin{quote}
The sound of the veena is very sweet.
\end{quote}

A long-necked string instrument of the south.

\begin{cousinweb}
\textbf{\kn{ವೀಣೆ}} (\emph{vīṇe}) is borrowed from Sanskrit \emph{vīṇā}, \textquotedblleft{}a lute\textquotedblright{} (CDIAL 12048); the final -ā became -\kn{ೆ}.
\end{cousinweb}

\begin{grammarlens}[title={What you've built}]
One more word for describing what you lived and felt.
\end{grammarlens}

\subsection*{Your turn}
\begin{itemize}
\raggedright
  \item \emph{Say it:} \emph{vīṇe}
  \item \emph{Say it:} \emph{vīṇe}, once more
  \item \emph{Say it:} \emph{vīṇeya nāda tumbā madhura.} — The sound of the veena is very sweet.
  \item \emph{Recall:} say the Kannada for to play an instrument, then the Kannada for a musical instrument, then say \emph{vīṇe} again
\end{itemize}

\begin{tcolorbox}[breakable,colback=teal!4,colframe=teal!35!black,title={Before you move on}]
What does \kn{ವೀಣೆ} mean? (\textquotedblleft{}A veena\textquotedblright{}.) Say it once more.
\end{tcolorbox}
\section[\texorpdfstring{\kn{ಕೊಳಲು} (koḷalu)}{ಕೊಳಲು (koḷalu)}]{\kn{ಕೊಳಲು} (koḷalu) — a flute}

\label{lesson:KA-C535-kolalu}



\begin{quote}
Before the new one: say the Kannada for a musical instrument, then the Kannada for a veena.
\end{quote}

\subsection*{You'll want to know: \kn{ಕೊಳಲು}}
\textbf{\kn{ಕೊಳಲು}} — \emph{koḷalu} — \textquotedblleft{}a flute\textquotedblright{}.

\begin{quote}
\kn{ದನ} \kn{ಕಾಯುವ} \kn{ಹುಡುಗ} \kn{ಕೊಳಲು} \kn{ನುಡಿಸುತ್ತಿದ್ದನು}.
\end{quote}

\begin{quote}
\emph{dana kāyuva huḍuga koḷalu nuḍisuttiddanu.}
\end{quote}

\begin{quote}
The boy minding the cattle was playing a flute.
\end{quote}

Made of bamboo, \kn{ಬಿದಿರು}.

\begin{cousinweb}
\textbf{\kn{ಕೊಳಲು}} (\emph{koḷalu}) is native Dravidian: DEDR 1818 lists the Kannada word for a flute, a fife, beside a word for a tube, a blow-pipe; \kn{ಕೊಳವೆ}, a pipe, is that second word.
\end{cousinweb}

\begin{grammarlens}[title={What you've built}]
One more word for describing what you lived and felt.
\end{grammarlens}

\subsection*{Your turn}
\begin{itemize}
\raggedright
  \item \emph{Say it:} \emph{koḷalu}
  \item \emph{Say it:} \emph{koḷalu}, once more
  \item \emph{Say it:} \emph{dana kāyuva huḍuga koḷalu nuḍisuttiddanu.} — The boy minding the cattle was playing a flute.
  \item \emph{Recall:} say the Kannada for a musical instrument, then the Kannada for a veena, then say \emph{koḷalu} again
\end{itemize}

\begin{tcolorbox}[breakable,colback=teal!4,colframe=teal!35!black,title={Before you move on}]
What does \kn{ಕೊಳಲು} mean? (\textquotedblleft{}A flute\textquotedblright{}.) Say it once more.
\end{tcolorbox}
\section[\texorpdfstring{\kn{ಪಿಟೀಲು} (piṭīlu)}{ಪಿಟೀಲು (piṭīlu)}]{\kn{ಪಿಟೀಲು} (piṭīlu) — a violin}

\label{lesson:KA-C535-pitilu}



\begin{quote}
Before the new one: say the Kannada for a veena, then the Kannada for a flute.
\end{quote}

\subsection*{You'll want to know: \kn{ಪಿಟೀಲು}}
\textbf{\kn{ಪಿಟೀಲು}} — \emph{piṭīlu} — \textquotedblleft{}a violin\textquotedblright{}.

\begin{quote}
\kn{ಸಂಗೀತ} \kn{ಕಚೇರಿಯಲ್ಲಿ} \kn{ಅವರು} \kn{ಪಿಟೀಲು} \kn{ನುಡಿಸಿದರು}.
\end{quote}

\begin{quote}
\emph{saṅgīta kacēriyalli avaru piṭīlu nuḍisidaru.}
\end{quote}

\begin{quote}
He played the violin at the concert.
\end{quote}

In southern concerts it follows the singer.

\begin{cousinweb}
\textbf{\kn{ಪಿಟೀಲು}} (\emph{piṭīlu}) is the Kannada name of the violin, which came to south India from Europe; the origin of the word is not settled.
\end{cousinweb}

\begin{grammarlens}[title={What you've built}]
One more word for describing what you lived and felt.
\end{grammarlens}

\subsection*{Your turn}
\begin{itemize}
\raggedright
  \item \emph{Say it:} \emph{piṭīlu}
  \item \emph{Say it:} \emph{piṭīlu}, once more
  \item \emph{Say it:} \emph{saṅgīta kacēriyalli avaru piṭīlu nuḍisidaru.} — He played the violin at the concert.
  \item \emph{Recall:} say the Kannada for a veena, then the Kannada for a flute, then say \emph{piṭīlu} again
\end{itemize}

\begin{tcolorbox}[breakable,colback=teal!4,colframe=teal!35!black,title={Before you move on}]
What does \kn{ಪಿಟೀಲು} mean? (\textquotedblleft{}A violin\textquotedblright{}.) Say it once more.
\end{tcolorbox}
\section[\texorpdfstring{\kn{ಬಾರಿಸು} (bārisu)}{ಬಾರಿಸು (bārisu)}]{\kn{ಬಾರಿಸು} (bārisu) — to beat a drum; to ring}

\label{lesson:KA-C535-barisu}



\begin{quote}
Before the new one: say the Kannada for a flute, then the Kannada for a violin.
\end{quote}

\subsection*{You'll want to know: \kn{ಬಾರಿಸು}}
\textbf{\kn{ಬಾರಿಸು}} — \emph{bārisu} — \textquotedblleft{}to beat a drum; to ring\textquotedblright{}.

\begin{quote}
\kn{ದೇವಸ್ಥಾನದಲ್ಲಿ} \kn{ಒಬ್ಬರು} \kn{ಗಂಟೆ} \kn{ಬಾರಿಸಿದರು}.
\end{quote}

\begin{quote}
\emph{dēvasthānadalli obbaru gaṇṭe bārisidaru.}
\end{quote}

\begin{quote}
Someone rang the bell at the temple.
\end{quote}

Past (I …): \kn{ಬಾರಿಸಿದೆ}, \emph{bāriside}.

Drums and bells take \kn{ಬಾರಿಸು}; strings and flutes take \kn{ನುಡಿಸು}.

\begin{cousinweb}
\textbf{\kn{ಬಾರಿಸು}} (\emph{bārisu}) is native Dravidian: DEDR 5380 lists Kannada \emph{bāṟisu}, \textquotedblleft{}to strike or beat with a whip\textquotedblright{}.
\end{cousinweb}

\begin{grammarlens}[title={What you've built}]
One more word for describing what you lived and felt.
\end{grammarlens}

\subsection*{Your turn}
\begin{itemize}
\raggedright
  \item \emph{Say it:} \emph{bārisu}
  \item \emph{Say it:} \emph{bārisu}, once more
  \item \emph{Say it:} \emph{dēvasthānadalli obbaru gaṇṭe bārisidaru.} — Someone rang the bell at the temple.
  \item \emph{Recall:} say the Kannada for a flute, then the Kannada for a violin, then say \emph{bārisu} again
\end{itemize}

\begin{tcolorbox}[breakable,colback=teal!4,colframe=teal!35!black,title={Before you move on}]
What does \kn{ಬಾರಿಸು} mean? (\textquotedblleft{}To beat a drum; to ring\textquotedblright{}.) Say it once more.
\end{tcolorbox}
